\lyseis{Α6.1}{Ιδιότητες της ισότητας και η έννοια της εξίσωσης}
\sol{1} α) Σωστό. \quad β) Λάθος: $x=7+3=10$. \quad γ) Σωστό. \quad δ) Σωστό (διαιρώ με $3$). \quad ε) Σωστό (πολλαπλασιάζω με $3$). \quad στ) Λάθος: $2\cdot6=12\neq6$, η λύση είναι το $3$.
\sol{2} $x+200=500$. Αφαιρώ $200$ και από τα δύο μέλη: $x=300$ g. Έλεγχος: $300+200=500$.
\sol{3} α) $\alpha=7-2=5$ \quad β) $\kappa=10+3=13$ \quad γ) $\lambda=3$ \quad δ) $\rho=45:9=5$.
\sol{4} α) Ναι: $6+3=9$. \quad β) Ναι: $9-2=7$. \quad γ) Όχι: $2\cdot20=40\neq10$ (η λύση είναι το $5$). \quad δ) Όχι: $12-22=-10\neq10$ (η λύση είναι το $2$). \quad ε) Ναι: $12:4=3$.
\sol{5} α) $33$ \quad β) $8$ \quad γ) $53$ \quad δ) $48$ \quad ε) $3$ \quad στ) $200$ \quad ζ) $2$.
\sol{6} α) $x=-5$, έλεγχος $-5\cdot(-5)=25$ \quad β) $x=-12+8=-4$, έλεγχος $-4-8=-12$ \quad γ) $x=3$, έλεγχος $-6\cdot3=-18$ \quad δ) $-x=11+8=19$, $x=-19$, έλεγχος $-8+19=11$.
\sol{7} α) $2{,}25$ \quad β) $x=\frac{9}{4}\cdot2=\frac{9}{2}$ \quad γ) $x=\frac{5}{6}+\frac{1}{3}=\frac{7}{6}$ \quad δ) $x=10\cdot\frac{3}{5}=6$.
\sol{8} α) $x+17=52$, $x=35$ \quad β) $3x=84$, $x=28$ \quad γ) $50-x=18$, $x=32$ \quad δ) $x:2=7$, $x=14$.
\sol{9} $2x=1\,200$, άρα $x=600$ g. Έλεγχος: $2\cdot600=1\,200$.
\sol{10} α) $4x=54$, $x=13{,}5$ € \quad β) $36-x=11{,}5$, $x=24{,}5$ €.
\sol{11} α) $4x=98$, $x=24{,}5$ cm. \quad β) $3x=4{,}5$, $x=1{,}5$ m.
\sol{12} α) Ταυτότητα (αντιμεταθετική ιδιότητα). \quad β) Αδύνατη. \quad γ) Αδύνατη: το αριστερό μέλος είναι πάντα κατά $4$ μικρότερο. \quad δ) Ταυτότητα: $5x-5x=0$ και $0=0$.
\sol{13} Κάθε αριθμός επί $0$ κάνει $0$, άρα η $0\cdot x=0$ ισχύει για κάθε $x$. Κανένας αριθμός επί $0$ δεν κάνει $5$, άρα η $0\cdot x=5$ δεν έχει λύση.
\sol{14} Στην πρώτη πρόσθεσε αντί να αφαιρέσει: από την $\alpha-x=\beta$ προκύπτει $x=\alpha-\beta$, δηλαδή $x=12-5=7$ (έλεγχος $12-7=5$). Στη δεύτερη πρέπει $x=24:6=4$ (έλεγχος $24:4=6$).
\sol{15} Ελεύθερη απάντηση. Π.χ. $x+4=10$ (λύση $6$), $x+1=x+2$ (αδύνατη), $x+3=3+x$ (ταυτότητα).
\sol{16} α) Ισοδύναμες: και οι δύο έχουν μόνη λύση το $4$. \quad β) Όχι: η πρώτη έχει λύση το $3$ και η δεύτερη το $1$. \quad γ) Ισοδύναμες: και οι δύο είναι αδύνατες, δηλαδή έχουν τις ίδιες (καμία) λύσεις.
\sol{17} $3x=2{,}4$, άρα $x=0{,}8$.
\sol{18} Ελεύθερη απάντηση. Π.χ. «Η Ζωή είχε $x$ ευρώ, πήρε άλλα $12$ και έχει τώρα $30$»: $x=30-12=18$.
